\documentclass[12pt,a4paper]{article}
\usepackage[utf8]{vietnam}
\usepackage{amsmath,amssymb}
\usepackage{tabularx}
\usepackage{tikz}
\usepackage{xcolor}
\usepackage[top=0.6cm, bottom=0.6cm, left=0.6cm, right=0.6cm]{geometry}
\usepackage[dethi]{ex_test}

\SAOPTN{widthSA=2.6,heightSA=0.62}
\setlength{\parskip}{0pt}

% Dong ke cham mo chuan viet tay
\newcommand{\dongke}[1]{%
  \begingroup%
  \setlength{\parskip}{0pt}%
  \baselineskip=18.5pt%
  \vspace{1.0mm}%
  \foreach \i in {1,...,#1}{%
    \noindent{\color{gray!60}\dotfill}\par%
  }%
  \vspace{0.6mm}%
  \endgroup%
}

% Dinh nghia o bai tap tu luan trong bang
\newcommand{\oBaiTap}[3][6]{%
  \vspace{1.0mm}%
  \noindent\textbf{#2} #3\par%
  \dongke{#1}%
}

\begin{document}
\pagestyle{empty}
% =============================================================================
% ========================= MA DE 101 - MAT TRUOC (TRANG 1/2) =============
% =============================================================================
\noindent
\begin{tabularx}{\linewidth}{|X|c|}
\hline
\textbf{PHIẾU KIỂM TRA ĐỊNH KỲ TOÁN 7 - TUẦN 07 - BUỔI 01} & \textbf{Mã đề: 101 - Trang 1/2} \\
\hline
\end{tabularx}

\vspace{1.2mm}

\noindent
\begin{tabularx}{\linewidth}{|p{3.2cm}|X|}
\hline
\centering\textbf{Điểm số} & \centering\textbf{Lời nhận xét của thầy cô} \tabularnewline
\rule{0pt}{1.8cm} & \\
\hline
\end{tabularx}

\vspace{1.2mm}

\noindent\textbf{Câu 1 (0,25 điểm).} Giá trị của $\sqrt{144}$ bằng bao nhiêu? \shortans{$12$}\par\vspace{0.3mm}
\noindent\textbf{Câu 2 (0,25 điểm).} Tính giá trị của biểu thức $\lvert -5{,}2\rvert - \lvert 3{,}2\rvert$. \shortans{$2$}\par\vspace{0.3mm}
\noindent\textbf{Câu 3 (0,25 điểm).} Làm tròn số $28{,}647$ đến hàng phần mười. \shortans{$28{,}6$}\par\vspace{0.3mm}
\noindent\textbf{Câu 4 (0,25 điểm).} Cho hai góc $\widehat{xOy}$ và $\widehat{yOz}$ kề bù, biết $\widehat{xOy}=75^\circ$. Tính số đo $\widehat{yOz}$. \shortans{$105^\circ$}\par\vspace{0.3mm}
\noindent\textbf{Câu 5 (0,25 điểm).} Hai đường thẳng cắt nhau tạo thành một góc có số đo $48^\circ$. Góc đối đỉnh với góc đó có số đo bằng bao nhiêu? \shortans{$48^\circ$}\par\vspace{0.3mm}
\noindent\textbf{Câu 6 (0,25 điểm).} Cho tia $Ot$ là tia phân giác của góc $\widehat{mOn}=110^\circ$. Tính số đo $\widehat{mOt}$. \shortans{$55^\circ$}\par\vspace{0.3mm}
\noindent\textbf{Câu 7 (0,25 điểm).} Tính giá trị của biểu thức $\left(-\frac{1}{3}\right)^2+\frac{8}{9}$. \shortans{$1$}\par\vspace{0.3mm}
\noindent\textbf{Câu 8 (0,25 điểm).} Qua điểm $M$ nằm ngoài đường thẳng $a$, có bao nhiêu đường thẳng song song với $a$? \shortans{$1$}\par\vspace{0.3mm}
\noindent\textbf{Câu 9 (0,25 điểm).} Tìm số thực $x\geq 0$, biết $\sqrt{x}=9$. \shortans{$81$}\par\vspace{0.3mm}
\noindent\textbf{Câu 10 (0,25 điểm).} Một cuốn sách có giá $80\,000$ đồng được giảm giá $15\%$. Số tiền được giảm là bao nhiêu đồng? \shortans{$12\,000$ đồng}\par\vspace{0.3mm}

\vspace{0.5mm}
\noindent\textbf{Câu 11 (2,0 điểm). Thực hiện phép tính (hợp lý nếu có thể):}
\vspace{0.5mm}

\noindent
\begin{tabularx}{\linewidth}{|X|X|}
\hline
\oBaiTap[10]{a)}{$\frac{5}{14}+\left(-\frac{7}{3}\right)+\frac{9}{14}$} & \oBaiTap[10]{b)}{$\frac{7}{11}\cdot\left(-\frac{8}{15}\right)+\frac{7}{11}\cdot\left(-\frac{7}{15}\right)$} \\
\hline
\oBaiTap[10]{c)}{$\left(-\frac{1}{2}\right)^2\cdot 12-2{,}5\cdot 4$} & \oBaiTap[10]{d)}{$\sqrt{\frac{36}{25}}-\lvert -0{,}2\rvert+\sqrt{49}$} \\
\hline
\end{tabularx}

\newpage
% =============================================================================
% ========================= MA DE 101 - MAT SAU (TRANG 2/2) ==============
% =============================================================================
\noindent
\begin{tabularx}{\linewidth}{|X|}
\hline
\centering\textbf{TOÁN 7 - TUẦN 07 - BUỔI 01 | Mã đề: 101 - Trang 2/2} \\
\hline
\end{tabularx}

\vspace{1.2mm}

\noindent\textbf{Câu 12 (2,0 điểm). Tìm $x$, biết:}
\vspace{0.5mm}

\noindent
\begin{tabularx}{\linewidth}{|X|X|}
\hline
\oBaiTap[5]{a)}{$x-\frac{1}{4}=\frac{5}{12}$} & \oBaiTap[5]{b)}{$\frac{3}{5}x+\frac{1}{2}=\frac{11}{10}$} \\
\hline
\oBaiTap[5]{c)}{$\left(x+\frac{1}{3}\right)^2=\frac{16}{25}$} & \oBaiTap[5]{d)}{$\lvert 3x-1\rvert=5$} \\
\hline
\end{tabularx}

\vspace{1.2mm}

\noindent
\begin{tabularx}{\linewidth}{|X|}
\hline
\vspace{0.5mm}
\textbf{Câu 13 (2,5 điểm).} Cho hình vẽ, đường thẳng $c$ cắt hai đường thẳng $a$ và $b$ lần lượt tại $A$ và $B$. Các góc tại mỗi giao điểm được đánh số $1,2,3,4$ theo cùng một chiều, trong đó $\widehat{A}_1$ và $\widehat{B}_1$ ở vị trí đồng vị. Biết $\widehat{A}_1=\widehat{B}_1=60^\circ$.\par
\textbf{a) (1,25 điểm)} Chứng minh rằng $a\parallel b$.\par
\textbf{b) (1,25 điểm)} Tính số đo các góc $\widehat{A}_2,\widehat{A}_3,\widehat{A}_4$ và $\widehat{B}_2,\widehat{B}_3,\widehat{B}_4$.\par
\vspace{0.4mm}
\noindent\textit{Bài giải:}\dongke{11}
\\ \hline
\end{tabularx}

\vspace{1.2mm}

\noindent
\begin{tabularx}{\linewidth}{|X|}
\hline
\vspace{0.5mm}
\textbf{Câu 14 (1,0 điểm).} Cho biểu thức $A=\frac{1}{2}+\frac{1}{2^2}+\frac{1}{2^3}+\dots+\frac{1}{2^{2026}}$. Hãy so sánh $A$ với $1$.\par
\vspace{0.4mm}
\noindent\textit{Bài giải:}\dongke{8}
\\ \hline
\end{tabularx}

\newpage
% =============================================================================
% ========================= MA DE 102 - MAT TRUOC (TRANG 1/2) =============
% =============================================================================
\noindent
\begin{tabularx}{\linewidth}{|X|c|}
\hline
\textbf{PHIẾU KIỂM TRA ĐỊNH KỲ TOÁN 7 - TUẦN 07 - BUỔI 01} & \textbf{Mã đề: 102 - Trang 1/2} \\
\hline
\end{tabularx}

\vspace{1.2mm}

\noindent
\begin{tabularx}{\linewidth}{|p{3.2cm}|X|}
\hline
\centering\textbf{Điểm số} & \centering\textbf{Lời nhận xét của thầy cô} \tabularnewline
\rule{0pt}{1.8cm} & \\
\hline
\end{tabularx}

\vspace{1.2mm}

\noindent\textbf{Câu 1 (0,25 điểm).} Giá trị của $\sqrt{169}$ bằng bao nhiêu? \shortans{$13$}\par\vspace{0.3mm}
\noindent\textbf{Câu 2 (0,25 điểm).} Tính giá trị của biểu thức $\lvert -6{,}4\rvert - \lvert 3{,}4\rvert$. \shortans{$3$}\par\vspace{0.3mm}
\noindent\textbf{Câu 3 (0,25 điểm).} Làm tròn số $37{,}482$ đến hàng phần mười. \shortans{$37{,}5$}\par\vspace{0.3mm}
\noindent\textbf{Câu 4 (0,25 điểm).} Cho hai góc $\widehat{xOy}$ và $\widehat{yOz}$ kề bù, biết $\widehat{xOy}=55^\circ$. Tính số đo $\widehat{yOz}$. \shortans{$125^\circ$}\par\vspace{0.3mm}
\noindent\textbf{Câu 5 (0,25 điểm).} Hai đường thẳng cắt nhau tạo thành một góc có số đo $62^\circ$. Góc đối đỉnh với góc đó có số đo bằng bao nhiêu? \shortans{$62^\circ$}\par\vspace{0.3mm}
\noindent\textbf{Câu 6 (0,25 điểm).} Cho tia $Ot$ là tia phân giác của góc $\widehat{mOn}=120^\circ$. Tính số đo $\widehat{mOt}$. \shortans{$60^\circ$}\par\vspace{0.3mm}
\noindent\textbf{Câu 7 (0,25 điểm).} Tính giá trị của biểu thức $\left(-\frac{2}{5}\right)^2+\frac{21}{25}$. \shortans{$1$}\par\vspace{0.3mm}
\noindent\textbf{Câu 8 (0,25 điểm).} Hai đường thẳng phân biệt cùng song song với đường thẳng thứ ba thì có quan hệ gì với nhau? \shortans{Song song}\par\vspace{0.3mm}
\noindent\textbf{Câu 9 (0,25 điểm).} Tìm số thực $x\geq 0$, biết $\sqrt{x}=10$. \shortans{$100$}\par\vspace{0.3mm}
\noindent\textbf{Câu 10 (0,25 điểm).} Một chiếc ba lô có giá $240\,000$ đồng được giảm giá $15\%$. Số tiền được giảm là bao nhiêu đồng? \shortans{$36\,000$ đồng}\par\vspace{0.3mm}

\vspace{0.5mm}
\noindent\textbf{Câu 11 (2,0 điểm). Thực hiện phép tính (hợp lý nếu có thể):}
\vspace{0.5mm}

\noindent
\begin{tabularx}{\linewidth}{|X|X|}
\hline
\oBaiTap[10]{a)}{$\frac{4}{15}+\left(-\frac{5}{2}\right)+\frac{11}{15}$} & \oBaiTap[10]{b)}{$\frac{9}{13}\cdot\left(-\frac{5}{17}\right)+\frac{9}{13}\cdot\left(-\frac{12}{17}\right)$} \\
\hline
\oBaiTap[10]{c)}{$\left(-\frac{1}{3}\right)^2\cdot 18-3{,}5\cdot 2$} & \oBaiTap[10]{d)}{$\sqrt{\frac{64}{25}}-\lvert -0{,}6\rvert+\sqrt{64}$} \\
\hline
\end{tabularx}

\newpage
% =============================================================================
% ========================= MA DE 102 - MAT SAU (TRANG 2/2) ==============
% =============================================================================
\noindent
\begin{tabularx}{\linewidth}{|X|}
\hline
\centering\textbf{TOÁN 7 - TUẦN 07 - BUỔI 01 | Mã đề: 102 - Trang 2/2} \\
\hline
\end{tabularx}

\vspace{1.2mm}

\noindent\textbf{Câu 12 (2,0 điểm). Tìm $x$, biết:}
\vspace{0.5mm}

\noindent
\begin{tabularx}{\linewidth}{|X|X|}
\hline
\oBaiTap[5]{a)}{$x-\frac{1}{6}=\frac{7}{18}$} & \oBaiTap[5]{b)}{$\frac{4}{7}x+\frac{1}{3}=\frac{19}{21}$} \\
\hline
\oBaiTap[5]{c)}{$\left(x-\frac{1}{4}\right)^2=\frac{25}{16}$} & \oBaiTap[5]{d)}{$\lvert 3x+2\rvert=8$} \\
\hline
\end{tabularx}

\vspace{1.2mm}

\noindent
\begin{tabularx}{\linewidth}{|X|}
\hline
\vspace{0.5mm}
\textbf{Câu 13 (2,5 điểm).} Cho hình vẽ, đường thẳng $c$ cắt hai đường thẳng $a$ và $b$ lần lượt tại $A$ và $B$. Các góc tại mỗi giao điểm được đánh số $1,2,3,4$ theo cùng một chiều, trong đó $\widehat{A}_1$ và $\widehat{B}_1$ ở vị trí đồng vị. Biết $\widehat{A}_1=\widehat{B}_1=70^\circ$.\par
\textbf{a) (1,25 điểm)} Chứng minh rằng $a\parallel b$.\par
\textbf{b) (1,25 điểm)} Tính số đo các góc $\widehat{A}_2,\widehat{A}_3,\widehat{A}_4$ và $\widehat{B}_2,\widehat{B}_3,\widehat{B}_4$.\par
\vspace{0.4mm}
\noindent\textit{Bài giải:}\dongke{11}
\\ \hline
\end{tabularx}

\vspace{1.2mm}

\noindent
\begin{tabularx}{\linewidth}{|X|}
\hline
\vspace{0.5mm}
\textbf{Câu 14 (1,0 điểm).} Cho biểu thức $B=\frac{1}{3}+\frac{1}{3^2}+\frac{1}{3^3}+\dots+\frac{1}{3^{2026}}$. Hãy so sánh $B$ với $\frac{1}{2}$.\par
\vspace{0.4mm}
\noindent\textit{Bài giải:}\dongke{8}
\\ \hline
\end{tabularx}

\newpage
% =============================================================================
% ========================= MA DE 103 - MAT TRUOC (TRANG 1/2) =============
% =============================================================================
\noindent
\begin{tabularx}{\linewidth}{|X|c|}
\hline
\textbf{PHIẾU KIỂM TRA ĐỊNH KỲ TOÁN 7 - TUẦN 07 - BUỔI 01} & \textbf{Mã đề: 103 - Trang 1/2} \\
\hline
\end{tabularx}

\vspace{1.2mm}

\noindent
\begin{tabularx}{\linewidth}{|p{3.2cm}|X|}
\hline
\centering\textbf{Điểm số} & \centering\textbf{Lời nhận xét của thầy cô} \tabularnewline
\rule{0pt}{1.8cm} & \\
\hline
\end{tabularx}

\vspace{1.2mm}

\noindent\textbf{Câu 1 (0,25 điểm).} Giá trị của $\sqrt{196}$ bằng bao nhiêu? \shortans{$14$}\par\vspace{0.3mm}
\noindent\textbf{Câu 2 (0,25 điểm).} Tính giá trị của biểu thức $\lvert -7{,}6\rvert - \lvert 3{,}6\rvert$. \shortans{$4$}\par\vspace{0.3mm}
\noindent\textbf{Câu 3 (0,25 điểm).} Làm tròn số $46{,}359$ đến hàng phần mười. \shortans{$46{,}4$}\par\vspace{0.3mm}
\noindent\textbf{Câu 4 (0,25 điểm).} Cho hai góc $\widehat{xOy}$ và $\widehat{yOz}$ kề bù, biết $\widehat{xOy}=65^\circ$. Tính số đo $\widehat{yOz}$. \shortans{$115^\circ$}\par\vspace{0.3mm}
\noindent\textbf{Câu 5 (0,25 điểm).} Hai đường thẳng cắt nhau tạo thành một góc có số đo $54^\circ$. Góc đối đỉnh với góc đó có số đo bằng bao nhiêu? \shortans{$54^\circ$}\par\vspace{0.3mm}
\noindent\textbf{Câu 6 (0,25 điểm).} Cho tia $Ot$ là tia phân giác của góc $\widehat{mOn}=130^\circ$. Tính số đo $\widehat{mOt}$. \shortans{$65^\circ$}\par\vspace{0.3mm}
\noindent\textbf{Câu 7 (0,25 điểm).} Tính giá trị của biểu thức $\left(-\frac{3}{5}\right)^2+\frac{16}{25}$. \shortans{$1$}\par\vspace{0.3mm}
\noindent\textbf{Câu 8 (0,25 điểm).} Hai đường thẳng phân biệt cùng vuông góc với đường thẳng thứ ba thì có quan hệ gì với nhau? \shortans{Song song}\par\vspace{0.3mm}
\noindent\textbf{Câu 9 (0,25 điểm).} Tìm số thực $x\geq 0$, biết $\sqrt{x}=11$. \shortans{$121$}\par\vspace{0.3mm}
\noindent\textbf{Câu 10 (0,25 điểm).} Một đôi dép có giá $160\,000$ đồng được giảm giá $15\%$. Số tiền được giảm là bao nhiêu đồng? \shortans{$24\,000$ đồng}\par\vspace{0.3mm}

\vspace{0.5mm}
\noindent\textbf{Câu 11 (2,0 điểm). Thực hiện phép tính (hợp lý nếu có thể):}
\vspace{0.5mm}

\noindent
\begin{tabularx}{\linewidth}{|X|X|}
\hline
\oBaiTap[10]{a)}{$\frac{7}{18}+\left(-\frac{8}{3}\right)+\frac{11}{18}$} & \oBaiTap[10]{b)}{$\frac{5}{12}\cdot\left(-\frac{11}{19}\right)+\frac{5}{12}\cdot\left(-\frac{8}{19}\right)$} \\
\hline
\oBaiTap[10]{c)}{$\left(-\frac{1}{2}\right)^3\cdot 16+1{,}5\cdot(-4)$} & \oBaiTap[10]{d)}{$\sqrt{\frac{49}{16}}-\lvert -0{,}75\rvert+\sqrt{81}$} \\
\hline
\end{tabularx}

\newpage
% =============================================================================
% ========================= MA DE 103 - MAT SAU (TRANG 2/2) ==============
% =============================================================================
\noindent
\begin{tabularx}{\linewidth}{|X|}
\hline
\centering\textbf{TOÁN 7 - TUẦN 07 - BUỔI 01 | Mã đề: 103 - Trang 2/2} \\
\hline
\end{tabularx}

\vspace{1.2mm}

\noindent\textbf{Câu 12 (2,0 điểm). Tìm $x$, biết:}
\vspace{0.5mm}

\noindent
\begin{tabularx}{\linewidth}{|X|X|}
\hline
\oBaiTap[5]{a)}{$x-\frac{1}{5}=\frac{7}{20}$} & \oBaiTap[5]{b)}{$\frac{5}{8}x-\frac{1}{4}=\frac{3}{8}$} \\
\hline
\oBaiTap[5]{c)}{$\left(x+\frac{1}{2}\right)^2=\frac{9}{4}$} & \oBaiTap[5]{d)}{$\lvert 2x-3\rvert=7$} \\
\hline
\end{tabularx}

\vspace{1.2mm}

\noindent
\begin{tabularx}{\linewidth}{|X|}
\hline
\vspace{0.5mm}
\textbf{Câu 13 (2,5 điểm).} Cho hình vẽ, đường thẳng $c$ cắt hai đường thẳng $a$ và $b$ lần lượt tại $A$ và $B$. Các góc tại mỗi giao điểm được đánh số $1,2,3,4$ theo cùng một chiều, trong đó $\widehat{A}_1$ và $\widehat{B}_1$ ở vị trí đồng vị. Biết $\widehat{A}_1=\widehat{B}_1=50^\circ$.\par
\textbf{a) (1,25 điểm)} Chứng minh rằng $a\parallel b$.\par
\textbf{b) (1,25 điểm)} Tính số đo các góc $\widehat{A}_2,\widehat{A}_3,\widehat{A}_4$ và $\widehat{B}_2,\widehat{B}_3,\widehat{B}_4$.\par
\vspace{0.4mm}
\noindent\textit{Bài giải:}\dongke{11}
\\ \hline
\end{tabularx}

\vspace{1.2mm}

\noindent
\begin{tabularx}{\linewidth}{|X|}
\hline
\vspace{0.5mm}
\textbf{Câu 14 (1,0 điểm).} Cho biểu thức $C=\frac{1}{4}+\frac{1}{4^2}+\frac{1}{4^3}+\dots+\frac{1}{4^{2026}}$. Hãy so sánh $C$ với $\frac{1}{3}$.\par
\vspace{0.4mm}
\noindent\textit{Bài giải:}\dongke{8}
\\ \hline
\end{tabularx}

\newpage
% =============================================================================
% ========================= MA DE 104 - MAT TRUOC (TRANG 1/2) =============
% =============================================================================
\noindent
\begin{tabularx}{\linewidth}{|X|c|}
\hline
\textbf{PHIẾU KIỂM TRA ĐỊNH KỲ TOÁN 7 - TUẦN 07 - BUỔI 01} & \textbf{Mã đề: 104 - Trang 1/2} \\
\hline
\end{tabularx}

\vspace{1.2mm}

\noindent
\begin{tabularx}{\linewidth}{|p{3.2cm}|X|}
\hline
\centering\textbf{Điểm số} & \centering\textbf{Lời nhận xét của thầy cô} \tabularnewline
\rule{0pt}{1.8cm} & \\
\hline
\end{tabularx}

\vspace{1.2mm}

\noindent\textbf{Câu 1 (0,25 điểm).} Giá trị của $\sqrt{225}$ bằng bao nhiêu? \shortans{$15$}\par\vspace{0.3mm}
\noindent\textbf{Câu 2 (0,25 điểm).} Tính giá trị của biểu thức $\lvert -8{,}8\rvert - \lvert 3{,}8\rvert$. \shortans{$5$}\par\vspace{0.3mm}
\noindent\textbf{Câu 3 (0,25 điểm).} Làm tròn số $54{,}728$ đến hàng phần mười. \shortans{$54{,}7$}\par\vspace{0.3mm}
\noindent\textbf{Câu 4 (0,25 điểm).} Cho hai góc $\widehat{xOy}$ và $\widehat{yOz}$ kề bù, biết $\widehat{xOy}=85^\circ$. Tính số đo $\widehat{yOz}$. \shortans{$95^\circ$}\par\vspace{0.3mm}
\noindent\textbf{Câu 5 (0,25 điểm).} Hai đường thẳng cắt nhau tạo thành một góc có số đo $76^\circ$. Góc đối đỉnh với góc đó có số đo bằng bao nhiêu? \shortans{$76^\circ$}\par\vspace{0.3mm}
\noindent\textbf{Câu 6 (0,25 điểm).} Cho tia $Ot$ là tia phân giác của góc $\widehat{mOn}=140^\circ$. Tính số đo $\widehat{mOt}$. \shortans{$70^\circ$}\par\vspace{0.3mm}
\noindent\textbf{Câu 7 (0,25 điểm).} Tính giá trị của biểu thức $\left(-\frac{1}{4}\right)^2+\frac{15}{16}$. \shortans{$1$}\par\vspace{0.3mm}
\noindent\textbf{Câu 8 (0,25 điểm).} Cho $a\parallel b$ và đường thẳng $c\perp a$. Đường thẳng $c$ có quan hệ gì với đường thẳng $b$? \shortans{Vuông góc}\par\vspace{0.3mm}
\noindent\textbf{Câu 9 (0,25 điểm).} Tìm số thực $x\geq 0$, biết $\sqrt{x}=12$. \shortans{$144$}\par\vspace{0.3mm}
\noindent\textbf{Câu 10 (0,25 điểm).} Một chiếc áo có giá $120\,000$ đồng được giảm giá $15\%$. Số tiền được giảm là bao nhiêu đồng? \shortans{$18\,000$ đồng}\par\vspace{0.3mm}

\vspace{0.5mm}
\noindent\textbf{Câu 11 (2,0 điểm). Thực hiện phép tính (hợp lý nếu có thể):}
\vspace{0.5mm}

\noindent
\begin{tabularx}{\linewidth}{|X|X|}
\hline
\oBaiTap[10]{a)}{$\frac{5}{16}+\left(-\frac{7}{2}\right)+\frac{11}{16}$} & \oBaiTap[10]{b)}{$\frac{6}{17}\cdot\left(-\frac{13}{21}\right)+\frac{6}{17}\cdot\left(-\frac{8}{21}\right)$} \\
\hline
\oBaiTap[10]{c)}{$\left(-\frac{2}{3}\right)^2\cdot 9-4{,}5\cdot 2$} & \oBaiTap[10]{d)}{$\sqrt{\frac{81}{16}}-\lvert -1{,}25\rvert+\sqrt{100}$} \\
\hline
\end{tabularx}

\newpage
% =============================================================================
% ========================= MA DE 104 - MAT SAU (TRANG 2/2) ==============
% =============================================================================
\noindent
\begin{tabularx}{\linewidth}{|X|}
\hline
\centering\textbf{TOÁN 7 - TUẦN 07 - BUỔI 01 | Mã đề: 104 - Trang 2/2} \\
\hline
\end{tabularx}

\vspace{1.2mm}

\noindent\textbf{Câu 12 (2,0 điểm). Tìm $x$, biết:}
\vspace{0.5mm}

\noindent
\begin{tabularx}{\linewidth}{|X|X|}
\hline
\oBaiTap[5]{a)}{$x-\frac{1}{3}=\frac{5}{12}$} & \oBaiTap[5]{b)}{$\frac{2}{5}x-\frac{1}{10}=\frac{3}{10}$} \\
\hline
\oBaiTap[5]{c)}{$\left(x-\frac{1}{5}\right)^2=\frac{16}{25}$} & \oBaiTap[5]{d)}{$\lvert 2x+5\rvert=9$} \\
\hline
\end{tabularx}

\vspace{1.2mm}

\noindent
\begin{tabularx}{\linewidth}{|X|}
\hline
\vspace{0.5mm}
\textbf{Câu 13 (2,5 điểm).} Cho hình vẽ, đường thẳng $c$ cắt hai đường thẳng $a$ và $b$ lần lượt tại $A$ và $B$. Các góc tại mỗi giao điểm được đánh số $1,2,3,4$ theo cùng một chiều, trong đó $\widehat{A}_1$ và $\widehat{B}_1$ ở vị trí đồng vị. Biết $\widehat{A}_1=\widehat{B}_1=65^\circ$.\par
\textbf{a) (1,25 điểm)} Chứng minh rằng $a\parallel b$.\par
\textbf{b) (1,25 điểm)} Tính số đo các góc $\widehat{A}_2,\widehat{A}_3,\widehat{A}_4$ và $\widehat{B}_2,\widehat{B}_3,\widehat{B}_4$.\par
\vspace{0.4mm}
\noindent\textit{Bài giải:}\dongke{11}
\\ \hline
\end{tabularx}

\vspace{1.2mm}

\noindent
\begin{tabularx}{\linewidth}{|X|}
\hline
\vspace{0.5mm}
\textbf{Câu 14 (1,0 điểm).} Cho biểu thức $D=\frac{1}{5}+\frac{1}{5^2}+\frac{1}{5^3}+\dots+\frac{1}{5^{2026}}$. Hãy so sánh $D$ với $\frac{1}{4}$.\par
\vspace{0.4mm}
\noindent\textit{Bài giải:}\dongke{8}
\\ \hline
\end{tabularx}

\newpage
\end{document}
